\begin{examplebox}
	\textbf{\textcolor{CExample}{例 1（基础）}}
	\textbf{\textcolor{CExample}{题目：}}
	设定义在 $\mathbb{R}$ 上的函数 $f(x)$ 以 $T=4$ 为周期，且有一条对称轴 $x=1$。写出 $f(x)$ 的所有对称轴方程。
	\textbf{\textcolor{CExample}{解题步骤：}}
	\begin{description}[style=nextline, leftmargin=2.8em, labelsep=0.5em, itemsep=0.45em]
		\item[\textbf{第一步（提取条件）：}]
			周期 $T=4$，已知对称轴 $a=1$。
			\par\textbf{理由：} 题目直接给出周期和对称轴数值。
		\item[\textbf{第二步（应用结论）：}]
			由结论一，全部对称轴为 $x = a + \dfrac{kT}{2}$，$k\in\mathbb{Z}$。
			\par\textbf{理由：} 该结论直接由周期与对称轴连锁性质得到。
		\item[\textbf{第三步（代入计算）：}]
			$x = 1 + \dfrac{k\cdot4}{2} = 1+2k$，$k\in\mathbb{Z}$。
			\par\textbf{理由：} 代入 $a=1$，$T=4$ 化简。
	\end{description}
	\textbf{\textcolor{CExample}{关键结论：}}
	\textbf{$x = 1+2k\ (k\in\mathbb{Z})$}。

	\medskip
	\textbf{\textcolor{CExample}{例 2（稍变形）}}
	\textbf{\textcolor{CExample}{题目：}}
	若函数 $f(x)$ 在 $\mathbb{R}$ 上以周期 $\pi$ 为周期，且关于点 $\left(\dfrac{\pi}{6},3\right)$ 中心对称，求 $f(x)$ 的所有对称中心坐标。
	\textbf{\textcolor{CExample}{解题步骤：}}
	\begin{description}[style=nextline, leftmargin=2.8em, labelsep=0.5em, itemsep=0.45em]
		\item[\textbf{第一步（提取条件）：}]
			周期 $T=\pi$，对称中心 $(b,c)=\left(\dfrac{\pi}{6},3\right)$。
			\par\textbf{理由：} 题目明确给出。
		\item[\textbf{第二步（应用结论二）：}]
			由结论二，所有对称中心为 $\left(b+\dfrac{kT}{2},\,c\right)$，$k\in\mathbb{Z}$。
			\par\textbf{理由：} 周期与对称中心连锁性质。
		\item[\textbf{第三步（代入数值）：}]
			$\left(\dfrac{\pi}{6}+\dfrac{k\pi}{2},\,3\right)$，$k\in\mathbb{Z}$。
			\par\textbf{理由：} 直接代入 $b=\pi/6$，$T=\pi$ 化简，注意纵坐标不变。
	\end{description}
	\textbf{\textcolor{CExample}{关键结论：}}
	\textbf{$\left(\dfrac{\pi}{6}+\dfrac{k\pi}{2},\,3\right)$（$k\in\mathbb{Z}$）}。

	\medskip
	\textbf{\textcolor{CExample}{例 3（综合应用）}}
	\textbf{\textcolor{CExample}{题目：}}
	已知定义在 $\mathbb{R}$ 上的函数 $f(x)$ 有两条不同的对称轴 $x=-1$ 和 $x=3$。求证 $f(x)$ 是周期函数，并求出它的一个正周期。
	\textbf{\textcolor{CExample}{解题步骤：}}
	\begin{description}[style=nextline, leftmargin=2.8em, labelsep=0.5em, itemsep=0.45em]
		\item[\textbf{第一步（列出对称轴条件）：}]
			由对称轴定义，有 $f(-1+x)=f(-1-x)$ 和 $f(3+x)=f(3-x)$ 对所有 $x$ 成立。
			\par\textbf{理由：} 对称轴条件等价于 $f(a+t)=f(a-t)$。
		\item[\textbf{第二步（应用推广结论）：}]
			由推广结论，周期 $T = 2|a_1-a_2| = 2\,|\,-1-3\,| = 8$。
			\par\textbf{理由：} 两条不同的对称轴决定周期为它们距离的两倍。
		\item[\textbf{第三步（验证周期性）：}]
			为严谨，也可直接推导：由 $f(t)=f(-2-t)$（从 $x=-1$ 得到）和 $f(t)=f(6-t)$（从 $x=3$ 得到），得 $f(-2-t)=f(6-t)$，令 $u=-2-t$，则 $f(u)=f(8+u)$，故 $8$ 是周期。
			\par\textbf{理由：} 代换法将对称轴关系转化为周期关系，验证结论正确。
	\end{description}
	\textbf{\textcolor{CExample}{关键结论：}}
	\textbf{$f(x+8)=f(x)$，周期 $T=8$}。
\end{examplebox}
